\chapter{Localisation and Dimension}

\begin{center}\itshape
The fraction $2/7$ belongs to $\Z_{(3)}$, but $1/3$ does not. We have allowed every denominator that avoids $3$. What information survives this choice? Exactly the information supported at the prime $3$. Localisation lets us make the same choice in any ring.
\end{center}

Suppose that a denominator has caused trouble in a calculation. Instead of ruling it out one fraction at a time, we can enlarge the ring once and make that denominator a unit. In $\Z[1/6]$, both $2$ and $3$ become invertible. In $\Z_{(3)}$, every prime except $3$ becomes invertible. These are different enlargements, chosen for different jobs.

The geometric version starts with a point $a$ of a curve. If a polynomial $g$ does not vanish at $a$, its reciprocal can be used near $a$. Fractions $f/g$ with $g(a)\ne0$ therefore record the functions available near that point. Wolfgang Krull helped make such local rings central to commutative algebra in the 1930s. The common operation behind these examples is \emph{localisation}.

Prime ideals also give a way to count dimensions. In $\C[x,y]$ the chain
\[
(0)\subsetneq(y)\subsetneq(x,y)
\]
records a surface, a curve in it, and a point on the curve. There are two strict inclusions. Krull's definition counts inclusions of prime ideals, so it applies even when a ring has no familiar geometric picture. Proving that this count gives $n$ for $K[x_1,\ldots,x_n]$ takes more work. Noether's normalisation lemma supplies the needed reduction.

There is a space behind every commutative ring: its prime spectrum. Grothendieck made this viewpoint part of the language of algebraic geometry. Mumford drew $\operatorname{Spec}\Z[x]$ as an arithmetic surface, with a fibre over each rational prime. We will justify the ideals in Figure~\ref{fig:specz}; the drawing is a guide to their inclusions, rather than an ordinary Euclidean surface.

All rings here are commutative and have $1$. We allow the zero ring, and flag the one theorem below that must exclude it. Ring homomorphisms preserve $1$.

\section{Rings of fractions}

To admit a denominator $s$, we must admit its powers. To admit $s$ and $t$, we must admit $st$. A subset $S\subseteq R$ is \emph{multiplicative} when $1\in S$ and it is closed under multiplication. Three choices will recur: $\{1,f,f^2,\ldots\}$; $R\setminus\mathfrak p$ for a prime ideal $\mathfrak p$; and the nonzero elements of a domain.

Ordinary cross-multiplication needs a small repair when there are zero divisors. We declare two fractions equal if an allowed denominator kills the difference obtained by cross-multiplication.

\begin{definition}\label{def:localisation}
Let $S\subseteq R$ be multiplicative. The \emph{localisation} $S^{-1}R$ consists of equivalence classes of pairs $(r,s)\in R\times S$, where
\[
(r,s)\sim(r',s')\quad\Longleftrightarrow\quad
 t(rs'-r's)=0\text{ for some }t\in S.
\]
Write the class as $r/s$, and put
\[
\frac rs+\frac{r'}{s'}=\frac{rs'+r's}{ss'},\qquad
\frac rs\frac{r'}{s'}=\frac{rr'}{ss'}.
\]
The canonical homomorphism is $\iota(r)=r/1$. We write $R_f$ when $S=\{f^n:n\ge0\}$, and $R_{\mathfrak p}$ when $S=R\setminus\mathfrak p$. For an $R$-module $M$, define $S^{-1}M$ by the same relation on pairs $(m,s)$, with the corresponding addition and scalar multiplication.
\end{definition}

Why retain the extra $t$? In $R=\Z/6$, localising at the powers of $2$ makes $2$ a unit. Since $2\cdot3=0$, the image of $3$ must then be zero. The fractions $3/1$ and $0/1$ must represent the same element, even though their cross-products in $R$ differ. The factor $t=2$ makes them equal.

Here is the main equivalence-relation check. Suppose
$t(rs'-r's)=0$ and $t'(r's''-r''s')=0$. Multiply the first equality by $t's''$, the second by $ts$, and add. The result is
\[
 tt's'(rs''-r''s)=0.
\]
The multiplier lies in $S$, which proves transitivity. Reflexivity uses $t=1$, and symmetry changes a sign. For addition and multiplication, replacing either input by an equivalent pair gives, after cross-multiplication, a sum of multiples of the old relations; multiplying by their witnesses in $S$ kills that sum. Thus the operations are well defined. Their ring laws follow by using a common denominator. The identity is $1/1$.

In a domain with $0\notin S$, cancellation removes $t$. We recover the usual fractions inside the fraction field. Thus $\Z_{(p)}$ consists of the $a/b$ with $p\nmid b$, and $\Z[1/6]$ consists of fractions whose reduced denominators have no primes other than $2$ and $3$. Localising a domain at $(0)$ gives its fraction field. If $0\in S$, every fraction is zero: choose $t=0$ in the relation. This includes the case where $S$ contains a nilpotent element.

The construction is useful because a map that already permits these divisions extends in exactly one way.

\begin{proposition}[universal property]\label{prop:locuniv}
For each $s\in S$, $\iota(s)$ is a unit. If $\varphi\colon R\to T$ is a ring homomorphism and every $\varphi(s)$, $s\in S$, is a unit of $T$, there is a unique homomorphism $\bar\varphi\colon S^{-1}R\to T$ with $\varphi=\bar\varphi\circ\iota$. Also
\[
\iota(r)=0\quad\Longleftrightarrow\quad sr=0\text{ for some }s\in S.
\]
\end{proposition}

\begin{proof}
The inverse of $s/1$ is $1/s$. Since $r/s=(r/1)(s/1)^{-1}$, any extension must satisfy
\[
\bar\varphi(r/s)=\varphi(r)\varphi(s)^{-1}.
\]
We check that this formula is independent of the pair. If $t(rs'-r's)=0$, its image under $\varphi$ is zero. Cancel the unit $\varphi(t)$, then divide by the units $\varphi(s)\varphi(s')$. This gives the same value for $r/s$ and $r'/s'$. The sum and product formulas show that the extension is a homomorphism, and its value at $r/1$ is $\varphi(r)$. Finally, the defining relation says that $r/1=0/1$ precisely when some $s\in S$ kills $r$.
\end{proof}

In Chapter~1's language, localisation is left adjoint to the inclusion of the $R$-algebras in which the chosen elements of $S$ become units. More concretely, any calculation with the allowed fractions can be transported to such an algebra by the displayed formula.

For modules, the zero test is just as useful:
$m/s=0$ if and only if $tm=0$ for some $t\in S$. It explains why localisation preserves exact sequences, although tensor products in general do not.

\begin{proposition}[localisation is exact]\label{prop:locexact}
\leavevmode
\begin{enumerate}[label=(\alph*)]
\item The formula $\frac rs\otimes m\mapsto\frac{rm}s$ gives a natural isomorphism $S^{-1}R\otimes_RM\cong S^{-1}M$.
\item If $M'\xrightarrow fM\xrightarrow gM''$ is exact, its localisation is exact, with maps $m/s\mapsto f(m)/s$ and $m/s\mapsto g(m)/s$.
\item The $R$-module $S^{-1}R$ is flat. In particular, $\Q$ is flat over $\Z$.
\end{enumerate}
\end{proposition}

\begin{proof}
For (a), multiplication $(r/s,m)\mapsto rm/s$ is $R$-bilinear and respects equality of fractions. Chapter~3's tensor universal property gives the indicated map. It is onto because $m/s$ is the image of $1/s\otimes m$.

Every tensor is a finite sum. Bring its denominators to the common denominator $s=\prod_i s_i$, then move the numerators across the tensor sign. It becomes $1/s\otimes m$ for one $m\in M$. If its image $m/s$ vanishes, choose $t\in S$ with $tm=0$. The tensor itself is then
\[
\frac1s\otimes m=\frac{t}{ts}\otimes m
 =\frac1{ts}\otimes tm=0.
\]
This proves injectivity. Applying a module homomorphism on either side sends this formula to the same formula, which proves naturality.

For (b), the composite of the localised maps is zero. Conversely, suppose that $m/s$ maps to zero. Some $t\in S$ satisfies $tg(m)=g(tm)=0$. Exactness before localisation gives $tm=f(m')$ for some $m'\in M'$. Hence
\[
\frac ms=\frac{tm}{ts}=\frac{f(m')}{ts},
\]
which lies in the localised image. This also handles an injective first map or a surjective last map by adjoining a zero module at that end. In (c), combine (a) and (b): tensoring with $S^{-1}R$ preserves injections. Taking $R=\Z$ and $S=\Z\setminus\{0\}$ gives $\Q$.
\end{proof}

Which primes survive? A prime meeting $S$ cannot survive, because a proper ideal cannot contain a unit. All other primes survive, and none are added.

\begin{theorem}[primes of a ring of fractions]\label{thm:locprimes}
Extension $\mathfrak p\mapsto S^{-1}\mathfrak p$ and contraction $\mathfrak q\mapsto\iota^{-1}(\mathfrak q)$ are inverse, inclusion-preserving bijections between the primes of $R$ disjoint from $S$ and the primes of $S^{-1}R$. Here $S^{-1}\mathfrak p=\{a/s:a\in\mathfrak p,\ s\in S\}$.
\end{theorem}

\begin{proof}
If $\mathfrak q$ is prime in $S^{-1}R$, its inverse image is prime. This follows directly by applying $\iota$ to a product. Its inverse image avoids $S$ since elements of $S$ become units. Extending that inverse image returns $\mathfrak q$: if $r/s\in\mathfrak q$, multiplication by $s/1$ gives $r/1\in\mathfrak q$; the reverse inclusion follows by multiplying $r/1$ by the unit $1/s$.

Now take a prime $\mathfrak p$ avoiding $S$. Its extension is an ideal. Suppose $r/1=a/s$ with $a\in\mathfrak p$. For some $t\in S$ we have $t(sr-a)=0$, so $tsr=ta\in\mathfrak p$. Neither $t$ nor $s$ belongs to $\mathfrak p$. Primeness therefore gives $r\in\mathfrak p$. We have proved that contraction of the extension is exactly $\mathfrak p$; in particular, the extension is proper.

If $(r/s)(r'/s')$ belongs to this extension, multiplication by $ss'/1$ gives $rr'/1$ in the extension. Contracting shows that $rr'\in\mathfrak p$, so $r\in\mathfrak p$ or $r'\in\mathfrak p$. One fraction therefore belongs to the extension. This proves primeness. Extension and contraction preserve inclusion from their definitions. If $0\in S$, both sets of primes are empty, so the assertion remains valid.
\end{proof}

\begin{corollary}\label{cor:localring}
For a prime $\mathfrak p$ of $R$, the ring $R_{\mathfrak p}$ has just one maximal ideal, $\mathfrak pR_{\mathfrak p}$. Its primes correspond to the primes of $R$ contained in $\mathfrak p$.
\end{corollary}

\begin{proof}
An ideal avoids $R\setminus\mathfrak p$ exactly when it is contained in $\mathfrak p$, so the prime correspondence gives the second assertion. The ideal $\mathfrak pR_{\mathfrak p}$ is proper. Every fraction $r/s$ outside it has $r\notin\mathfrak p$ and inverse $s/r$. Thus a proper ideal cannot contain any such fraction. Every proper ideal lies in $\mathfrak pR_{\mathfrak p}$, making it the unique maximal ideal.
\end{proof}

\section{Local rings and Nakayama's lemma}

A \emph{local ring} has exactly one maximal ideal, usually denoted $\mathfrak m$. The field $k=R/\mathfrak m$ is its \emph{residue field}. In a local ring, the units are exactly the elements outside $\mathfrak m$. To see the nontrivial direction, a nonunit generates a proper ideal; Zorn's lemma puts that ideal in a maximal ideal, which must be $\mathfrak m$. We use the form of Zorn's lemma proved in Volume~7, Theorem~1.19. In a Noetherian ring, Proposition~\ref{prop:noetherian} already gives the required maximal ideal.

A field is local with maximal ideal $0$. The ring $\Z_{(p)}$ is local with residue field $\mathbb F_p$. The formal power series ring $K[[x]]$ is local: a series with nonzero constant term has an inverse determined coefficient by coefficient, and the remaining series form the ideal $(x)$. Germs of smooth functions at a point give another example; the maximal ideal consists of germs vanishing at that point.

A finitely generated module can be tested over the much simpler field $k$. The reason is that a coefficient in $\mathfrak m$ cannot cancel the leading $1$ in $1-a$.

\begin{theorem}[Nakayama's lemma]\label{thm:nakayama}
Let $(R,\mathfrak m)$ be local. If $M$ is a finitely generated $R$-module and $\mathfrak mM=M$, then $M=0$.
\end{theorem}

\begin{proof}
If $M\ne0$, choose a generating list $m_1,\ldots,m_n$ of smallest possible length; then $n\ge1$. Since $m_n\in\mathfrak mM$, substituting expressions in this generating list gives
\[
 m_n=a_1m_1+\cdots+a_nm_n,\qquad a_i\in\mathfrak m.
\]
Rearrange:
\[
 (1-a_n)m_n=\sum_{i<n}a_im_i.
\]
The element $1-a_n$ cannot lie in $\mathfrak m$, since otherwise adding $a_n$ would put $1$ there. It is therefore a unit. Dividing expresses $m_n$ using the preceding generators. This contradicts minimality. For $n=1$, the sum on the right is zero, so the argument says $m_1=0$, again a contradiction.
\end{proof}

For instance, if a finite list of generators becomes zero after reduction modulo $\mathfrak m$, it was generating the zero module all along. The next consequence tells us exactly how many generators are needed.

\begin{corollary}\label{cor:mingens}
Let $M$ be finitely generated over the local ring $(R,\mathfrak m)$, with residue field $k$. A list $m_1,\ldots,m_n$ generates $M$ if and only if its images span $M/\mathfrak mM$ over $k$. Every minimal generating list therefore has $\dim_k(M/\mathfrak mM)$ elements.
\end{corollary}

\begin{proof}
A generating list certainly gives a spanning list in the quotient. Conversely, let $N=\sum_iRm_i$ and suppose the images span. Then $M=N+\mathfrak mM$. Thus $M/N=\mathfrak m(M/N)$, and this quotient is finitely generated. Nakayama gives $M/N=0$, or $M=N$.

If a generating list has linearly dependent images, one image is a linear combination of the others. Those other images still span, so the corresponding shorter list still generates $M$ by the first part. A minimal list therefore gives a basis. Conversely, a basis in the quotient cannot remain spanning when any member is removed. The empty list covers $M=0$.
\end{proof}

Finite generation cannot be dropped. Over $\Z_{(p)}$, the nonzero module $\Q$ satisfies $p\Q=\Q$. Its elements allow division by arbitrarily high powers of $p$, so no finite list generates it. Another useful local test needs no finite-generation assumption at all.

\begin{proposition}\label{prop:localzero}
For any $R$-module $M$, the following are equivalent: $M=0$; $M_{\mathfrak p}=0$ for every prime $\mathfrak p$; and $M_{\mathfrak m}=0$ for every maximal ideal $\mathfrak m$. A module homomorphism is injective, respectively surjective, if and only if all its localisations at maximal ideals have that property.
\end{proposition}

\begin{proof}
Only the implication from maximal ideals back to $M$ needs proof. Assume all $M_{\mathfrak m}$ vanish, and choose $x\in M$. The ideal
\[
\operatorname{Ann}(x)=\{r\in R:rx=0\}
\]
is the \emph{annihilator} of $x$. If it were proper, it would lie in some maximal ideal $\mathfrak m$. The equality $x/1=0$ in $M_{\mathfrak m}$ gives $sx=0$ for an $s\notin\mathfrak m$. This puts $s$ in $\operatorname{Ann}(x)\subseteq\mathfrak m$, a contradiction. Thus $1$ annihilates $x$, so $x=0$.

For a map $u$, apply this result to $\ker u$ or $\operatorname{coker}u$. Exactness of localisation identifies their localisations with the kernel or cokernel of $u_{\mathfrak m}$. If $R$ is the zero ring, its only unital module is zero, so the empty family of maximal ideals causes no exception.
\end{proof}

\section{The spectrum of a ring}

A maximal ideal records a field quotient, but nonmaximal primes record irreducible pieces as well. Keeping all primes makes localisation and the topology fit together.

\begin{definition}\label{def:spec}
The \emph{spectrum} $\operatorname{Spec}R$ is the set of prime ideals of $R$. Its \emph{Zariski topology} has closed sets
\[
 V(I)=\{\mathfrak p\in\operatorname{Spec}R:I\subseteq\mathfrak p\},
\]
where $I$ ranges over ideals. A ring homomorphism $\varphi\colon R\to T$ induces the continuous map $\operatorname{Spec}T\to\operatorname{Spec}R$ given by $\mathfrak q\mapsto\varphi^{-1}(\mathfrak q)$.
\end{definition}

The topology follows from three identities. We have $V(0)=\operatorname{Spec}R$, $V(R)=\varnothing$, and
\[
V(I)\cup V(J)=V(IJ),\qquad
\bigcap_\lambda V(I_\lambda)=V\Bigl(\sum_\lambda I_\lambda\Bigr).
\]
For the union identity, a prime containing $IJ$ must contain $I$ or $J$: otherwise choose an element outside the prime from each ideal and multiply. The intersection identity is simply the meaning of containing every $I_\lambda$. The inverse image of $V(I)$ under the induced map is $V(\varphi(I)T)$, which proves continuity. Composition reverses in the same way as inverse images. Thus $\operatorname{Spec}$ is a contravariant functor.

The closure of the single point $\mathfrak p$ is $V(\mathfrak p)$. In particular, a point is closed exactly when its prime is maximal. Over an algebraically closed field, the weak Nullstellensatz identifies maximal ideals of $K[x_1,\ldots,x_n]$ with the points of $K^n$. A nonmaximal prime adds the generic point of an irreducible closed set: its closure is the whole set it represents.

For $\Z$, the points are $(0)$ and $(p)$ for the rational primes $p$. Proper closed subsets are finite sets of the $(p)$, including the empty set. The closure of $(0)$ is the whole spectrum, so every nonempty open subset contains $(0)$. This is the meaning of a \emph{generic point}; it is a topological statement, not a claim that $(0)$ equals every point.

Passing from $\Z$ to $\Z[x]$ produces both arithmetic fibres and polynomial curves. The prime correspondence already classifies them.

\begin{theorem}[the primes of $\Z\lbrack x\rbrack$]\label{thm:speczx}
Every prime ideal of $\Z[x]$ is exactly one of the following types:
\begin{enumerate}[label=(\alph*)]
\item $(0)$;
\item $(f)$, where $f$ is irreducible in $\Z[x]$: either a rational prime, or a primitive polynomial irreducible in $\Q[x]$;
\item $(p,f)$, where $p$ is a rational prime and the reduction $\bar f$ is irreducible in $\mathbb F_p[x]$. These are maximal ideals.
\end{enumerate}
\end{theorem}

\begin{proof}
Gauss's theorem, Volume~3, Theorem~6.18, makes $\Z[x]$ a UFD. Therefore every ideal in (b) is prime. The quotient by $(p,f)$ is $\mathbb F_p[x]/(\bar f)$, a field, which verifies (c). The zero ideal is prime since $\Z[x]$ is a domain.

Conversely, contract a prime $P$ to $\Z$. The contraction is either $(0)$ or $(p)$. If it is $(0)$, localise at the nonzero integers. The resulting ring is the PID $\Q[x]$, whose primes are $(0)$ and $(g)$ with $g$ irreducible. The zero case contracts to zero. In the other case, multiply $g$ by a rational scalar to obtain a primitive $f\in\Z[x]$. Its contraction is
\[
 P=f\Q[x]\cap\Z[x]=f\Z[x].
\]
For the last equality, write a rational polynomial $h$ as $ch_0$ with $h_0$ primitive integral and $c\in\Q$. Gauss's lemma says $fh_0$ is primitive. If $fh=cfh_0$ has integral coefficients, the denominator of $c$ divides every coefficient of the primitive polynomial $fh_0$, and must be $1$. Thus $h\in\Z[x]$.

If $P\cap\Z=(p)$, pass to $\Z[x]/(p)=\mathbb F_p[x]$. This is a PID, so the prime $P/(p)$ is zero or generated by an irreducible $\bar f$. Its inverse image is $(p)$ or $(p,f)$, as claimed.
\end{proof}

\begin{figure}[ht]
\centering
\fig{ch04-speczx}
\caption{After Mumford: the spectrum of $\Z[x]$. The bottom line is $\operatorname{Spec}\Z$. Above each prime $p$ stands the fibre $\operatorname{Spec}\mathbb F_p[x]$, a vertical line made of the generic point $(p)$ and the
closed points $(p,f)$, one for each monic irreducible $\bar f$ modulo $p$, of which a few are drawn in grey. The curve is the closure of the point $(x^2+1)$: it meets the fibre at $p$ in two points
$(p,x\pm a)$ when $p\equiv1\pmod4$, in one point $(p,x^2+1)$ of degree $2$ when $p\equiv3\pmod4$, and has its two branches meeting at the single point $(2,x+1)$ of the fibre at $2$, where $x^2+1=(x+1)^2$.}
\label{fig:specz}
\end{figure}

The fibre over $(p)$ consists of the primes contracting to $(p)$: its generic point is $(p)$, and its closed points are $(p,f)$. Localising at a particular prime selects an even smaller set. Figure~\ref{fig:v10new-ch05-local} shows this selection for the simpler ring $\Z$.

\begin{figure}[ht]
\centering
\fig{v10-new-ch05-local}
\caption{Localisation at $(3)$ retains the primes $(0)$ and $(3)$ of $\Z$. Each other $(p)$ meets the denominators, since $p$ becomes a unit, and disappears. Lines show inclusion of primes, not distances between them.}
\label{fig:v10new-ch05-local}
\end{figure}

\section{Integral extensions}

The element $1/2$ satisfies $2x-1=0$ over $\Z$, but that does not make it an algebraic integer. The leading coefficient matters: an equation whose leading coefficient is $1$ lets us reduce every higher power using lower ones.

\begin{definition}\label{def:integral}
For rings $A\subseteq B$, an element $b\in B$ is \emph{integral over $A$} if it satisfies a monic polynomial in $A[x]$. The extension $B/A$ is integral when every element of $B$ is integral over $A$.
\end{definition}

For example, $\sqrt5$ is integral over $\Z$, since it satisfies $x^2-5=0$. Volume~13 will study such elements in number fields. We need only their algebra here. A monic equation is equivalent to a finite module of powers; the converse uses the adjugate matrix.

\begin{proposition}\label{prop:integralcrit}
For $b\in B$, the following conditions are equivalent:
(a) $b$ is integral over $A$;
(b) $A[b]$ is finitely generated as an $A$-module;
(c) some subring $C\subseteq B$ containing $A$ and $b$ is finitely generated as an $A$-module.
\end{proposition}

\begin{proof}
If $b^n+a_{n-1}b^{n-1}+\cdots+a_0=0$, it expresses $b^n$ using lower powers. Multiply by further powers of $b$ and use induction. Every element of $A[b]$ is an $A$-linear combination of $1,b,\ldots,b^{n-1}$, giving (a)$\Rightarrow$(b). The implication (b)$\Rightarrow$(c) takes $C=A[b]$.

For (c)$\Rightarrow$(a), choose generators $c_1,\ldots,c_n$ of $C$. Express multiplication by $b$ as
\[
 bc_i=\sum_j a_{ij}c_j,\qquad a_{ij}\in A.
\]
With $c$ the column of generators, $(bI-(a_{ij}))c=0$. Multiply by the adjugate. The determinant $D=\det(bI-(a_{ij}))$ kills every $c_i$, and hence kills their $A$-linear combination $1\in C$. Thus $D=0$. Expanding $D$ gives a monic polynomial in $b$ with coefficients in $A$. If $C$ is the zero ring, the same conclusion follows from the monic polynomial $x$ and its only element $b=0$.
\end{proof}

\begin{corollary}\label{cor:integralclosure}
The elements of $B$ integral over $A$ form a subring, called the \emph{integral closure} of $A$ in $B$. Integrality is transitive: if $C$ is integral over $B$ and $B$ is integral over $A$, then $C$ is integral over $A$.
\end{corollary}

\begin{proof}
If $b$ and $b'$ satisfy monic equations of degrees $r$ and $s$, the products $b^ib'^j$, with $0\le i<r$ and $0\le j<s$, generate $A[b,b']$ as an $A$-module. This ring contains $b+b'$, $b-b'$ and $bb'$, so the criterion makes each integral. The elements $0$ and $1$ are integral as well, proving the subring assertion.

For transitivity, choose $c\in C$ and a monic equation for it over $B$. Only finitely many coefficients $b_0,\ldots,b_{n-1}$ occur. Adjoin these to $A$ one at a time. At each step their monic equations show that the new ring is a finite module over the preceding ring. Products of the finite generating lists give a finite generating list over $A$. Finally adjoin $c$, using its monic equation. The resulting ring $A[b_0,\ldots,b_{n-1},c]$ is finite as an $A$-module, so the criterion makes $c$ integral over $A$.
\end{proof}

A domain is \emph{integrally closed}, or \emph{normal}, if the elements of its fraction field integral over it already belong to it. Every UFD has this property. Indeed, put $b=a/d$ in lowest terms into a monic equation and clear denominators. It follows that $d$ divides $a^n$, so coprimeness makes $d$ a unit. By contrast, $\Z[\sqrt5]$ omits the integral element $(1+\sqrt5)/2$ of its fraction field. Normal domains of dimension one will be the setting of Chapter~6.

An integral extension cannot turn a nonfield domain into a field. More precisely, a nonzero ideal upstairs always leaves a nonzero trace downstairs.

\begin{lemma}\label{lem:intfield}
If $A\subseteq B$ are domains and $B$ is integral over $A$, then $B$ is a field if and only if $A$ is a field. Every nonzero ideal of $B$ has nonzero intersection with $A$.
\end{lemma}

\begin{proof}
Choose $0\ne b\in B$ and a monic equation of least degree for $b$. Its constant term $a_0$ is nonzero: otherwise cancel the nonzero factor $b$ in the domain $B$ to obtain a shorter monic equation. Rearranging gives
\[
 0\ne a_0=-b(b^{n-1}+a_{n-1}b^{n-2}+\cdots+a_1)\in A\cap bB.
\]
Taking $b$ in any nonzero ideal proves the ideal assertion. If $A$ is a field, divide this equality by $a_0$ to write $b^{-1}$ as an element of $B$. Thus $B$ is a field.

Conversely, suppose $B$ is a field and $0\ne a\in A$. The element $a^{-1}\in B$ is integral over $A$. Multiply a monic equation of degree $m$ for it by $a^{m-1}$. The resulting equation expresses $a^{-1}$ as an $A$-linear combination of $1,a,\ldots,a^{m-1}$. Therefore $a^{-1}\in A$, making $A$ a field.
\end{proof}

These facts control whole chains of prime ideals. We will contract a prime upstairs by intersecting it with $A$. A prime above $\mathfrak p$ means one whose contraction is $\mathfrak p$.

\begin{theorem}[Cohen and Seidenberg, 1946]\label{thm:goingup}
Suppose $A\subseteq B$ and $B$ is integral over $A$.
\begin{enumerate}[label=(\alph*)]
\item \emph{Incomparability:} if primes $\mathfrak q\subseteq\mathfrak q'$ of $B$ contract to the same prime of $A$, they are equal.
\item \emph{Lying over:} each prime $\mathfrak p$ of $A$ is the contraction of a prime of $B$.
\item \emph{Going up:} if $\mathfrak p\subseteq\mathfrak p'$ are primes of $A$ and $\mathfrak q\cap A=\mathfrak p$, some prime $\mathfrak q'\supseteq\mathfrak q$ contracts to $\mathfrak p'$.
\end{enumerate}
\end{theorem}

\begin{proof}
For (a), pass to the domains $A/(\mathfrak q\cap A)\subseteq B/\mathfrak q$. The latter is integral over the former, since monic equations survive in quotients. The ideal $\mathfrak q'/\mathfrak q$ has zero contraction. Lemma~\ref{lem:intfield} says a nonzero ideal could not do so, hence $\mathfrak q'/\mathfrak q=0$.

For (b), put $S=A\setminus\mathfrak p$. The injection $A\to B$ remains injective under localisation, by exactness. Thus $A_{\mathfrak p}\subseteq S^{-1}B$, and this is a nonzero integral extension: for $b/s$, divide a monic equation for $b$ by its leading power of $s$. Choose a maximal ideal $\mathfrak n$ of $S^{-1}B$. The quotient field $(S^{-1}B)/\mathfrak n$ is integral over the domain
$A_{\mathfrak p}/(\mathfrak n\cap A_{\mathfrak p})$. By Lemma~\ref{lem:intfield}, the latter is a field too. Therefore $\mathfrak n\cap A_{\mathfrak p}$ is maximal, and must equal the unique maximal ideal $\mathfrak pA_{\mathfrak p}$. Contract $\mathfrak n$ to $B$ to obtain $\mathfrak q$. Contracting once more to $A$ gives $\mathfrak p$, by the localisation prime correspondence.

For (c), use the injection $A/\mathfrak p\subseteq B/\mathfrak q$. The extension is integral. Part (b) supplies a prime above $\mathfrak p'/\mathfrak p$; its inverse image in $B$ contains $\mathfrak q$ and contracts to $\mathfrak p'$. These arguments require no domain assumption on the original rings, because the quotients in which we use the lemma are domains.
\end{proof}

\section{Krull dimension and Noether normalisation}

The chain $(0)\subsetneq(x)$ in $K[x]$ has length $1$, whereas the single prime $(0)$ in a field gives a chain of length $0$. Length counts the strict steps, not the number of primes written.

\begin{definition}\label{def:krulldim}
The \emph{Krull dimension} of $R$ is the supremum of the integers $m$ for which there is a chain $\mathfrak p_0\subsetneq\cdots\subsetneq\mathfrak p_m$ of prime ideals. For a ring with no primes, we use the value $-\infty$.
\end{definition}

Fields have dimension $0$, and PIDs other than fields have dimension $1$. The classification of primes of $\Z[x]$ gives dimension $2$: there is a chain $(0)\subsetneq(x^2+1)\subsetneq(5,x-2)$, and no longer chain can occur. For the upper bound, two distinct primes of type (b) cannot contain one another. Above a fixed rational prime $p$, the only inclusions are $(p)\subsetneq(p,f)$. A prime $(f)$ contracting to $(0)$ cannot contain $(p)$. Thus any nonzero nonmaximal prime can be followed by at most one further strict step.

\begin{corollary}\label{cor:dimintegral}
An integral extension $A\subseteq B$ has $\dim B=\dim A$.
\end{corollary}

\begin{proof}
Contract any finite prime chain of $B$. Incomparability makes every contracted step strict, proving $\dim B\le\dim A$. Conversely, lying over lifts the first member of any finite chain in $A$, and going up lifts its successive members while preserving inclusion. Their distinct contractions keep the lifted steps strict. Taking suprema gives the reverse inequality, including infinite dimension. If there are no primes, lying over and contraction show that both spectra are empty.
\end{proof}

For polynomial rings, the lower bound on dimension is easy. To prove an upper bound, we want to replace a nonconstant polynomial relation by a monic relation in one variable. A substitution with widely separated powers prevents cancellation of the leading term.

\begin{lemma}[Noether's change of variables]\label{lem:noethervar}
Let $K$ be a field, $n\ge1$, and $f\in K[x_1,\ldots,x_n]$ nonconstant. For a sufficiently large integer $r$, substituting $x_i=y_i+x_n^{r^i}$ for $i<n$ gives
\[
 f=cx_n^N+\text{lower powers of }x_n,
\]
with $c\in K\setminus\{0\}$ and the remaining coefficients in $K[y_1,\ldots,y_{n-1}]$.
\end{lemma}

\begin{proof}
Choose $r$ greater than every exponent in every monomial of $f$. For a monomial $x^\alpha$, the substitution has highest power of $x_n$ equal to
\[
 e(\alpha)=\alpha_n+\alpha_1r+\cdots+\alpha_{n-1}r^{n-1},
\]
with coefficient $1$. Since all digits $\alpha_i$ are below $r$, distinct exponent vectors have distinct values of $e(\alpha)$. Choose the monomial of $f$ with largest value. Its leading term cannot cancel with any other monomial. Its original coefficient is the required nonzero constant $c$. Nonconstancy ensures $N>0$. For $n=1$, no substitution is needed.
\end{proof}

For example, substitute $x=y+z^3$ in $x^2+xz+1$. Its highest power of $z$ is $z^6$; the remaining terms have smaller degree in $z$. The polynomial becomes monic in $z$, even over a finite field. The argument uses distinct exponents, not a choice of sufficiently many field elements.

\begin{theorem}[Noether normalisation, 1926]\label{thm:noethernorm}
Let $A\ne0$ be a finitely generated algebra over a field $K$. There are algebraically independent elements $z_1,\ldots,z_d\in A$ such that $A$ is integral over the polynomial subring $K[z_1,\ldots,z_d]$.
\end{theorem}

\begin{proof}
Induct on the length $n$ of a generating list $a_1,\ldots,a_n$. A nonzero unital $K$-algebra contains a copy of $K$. If the generators are algebraically independent, they themselves are the required $z_i$. This includes $n=0$, when $A=K$ and $d=0$.

Otherwise choose a nonzero polynomial $f$ with $f(a_1,\ldots,a_n)=0$. It cannot be a nonzero constant, because $K\to A$ is injective. Apply the change-of-variables lemma, and put
\[
 b_i=a_i-a_n^{r^i}\quad(i<n),\qquad A'=K[b_1,\ldots,b_{n-1}].
\]
The substituted relation for $a_n$ has leading coefficient $c\in K^\times$. Divide by $c$ to make it monic. Thus $a_n$ is integral over $A'$. Since $a_i=b_i+a_n^{r^i}$, we have $A=A'[a_n]$, an integral extension. The nonzero algebra $A'$ has fewer generators, so the induction hypothesis finds algebraically independent $z_1,\ldots,z_d$ over which it is integral. Transitivity of integrality gives the result for $A$.
\end{proof}

The exclusion $A\ne0$ matters: the zero ring cannot contain a copy even of $K$, the polynomial ring on an empty list. When $d=0$, the theorem simply says that $A$ is a finite-dimensional $K$-algebra. In general $A$ is a finite module over $K[z_1,\ldots,z_d]$, because it has finitely many algebra generators and all are integral.

On spectra, the inclusion gives a surjective map to affine $d$-space by lying over. Its fibres have finitely many points: after fixing a prime downstairs and passing to its residue field, the fibre algebra is finite-dimensional, so its prime quotients are fields and it has only finitely many maximal ideals. For the last finiteness assertion, any $r$ distinct maximal ideals give, by the Chinese remainder theorem, a quotient whose vector-space dimension is at least $r$. Over algebraically closed $K$, this also describes a surjection on ordinary affine points. Over an arbitrary field, surjectivity on $K$-points need not hold; $\R\subset\C$ is already a warning.

\begin{theorem}\label{thm:dimpoly}
For a field $K$ and $n\ge0$, $\dim K[x_1,\ldots,x_n]=n$. Consequently, if a finitely generated $K$-algebra $A$ is integral over a polynomial subring $K[z_1,\ldots,z_d]$ in algebraically independent elements, then $\dim A=d$.
\end{theorem}

\begin{proof}
Each ideal in
\[
(0)\subsetneq(x_1)\subsetneq(x_1,x_2)\subsetneq\cdots
 \subsetneq(x_1,\ldots,x_n)
\]
is prime, since its quotient is a polynomial ring over a field. Hence the dimension is at least $n$.

For the upper bound, induct on $n$. The case $n=0$ is the field $K$. Take a prime chain of length $m\ge1$ in $K[x_1,\ldots,x_n]$. We may replace its first member by $(0)$; the next member remains nonzero and the length is unchanged. Choose $0\ne f\in\mathfrak p_1$. It is nonconstant since $\mathfrak p_1$ is proper. Use Noether's substitution, an automorphism with inverse $y_i=x_i-x_n^{r^i}$, and divide by the nonzero constant leading coefficient. We may view $f$ as monic of positive degree in $x_n$ over $A'=K[y_1,\ldots,y_{n-1}]$.

Then $B=K[x_1,\ldots,x_n]/(f)$ is integral over the image of $A'$. That image is a copy of $A'$: a nonzero multiple of a monic polynomial of positive $x_n$-degree has positive $x_n$-degree, and cannot lie in $A'$. The integral-dimension corollary and induction give $\dim B=n-1$. The quotient chain
\[
\mathfrak p_1/(f)\subsetneq\cdots\subsetneq\mathfrak p_m/(f)
\]
has length $m-1$, so $m-1\le n-1$, and $m\le n$. A length-zero chain also obeys the bound. The final assertion now follows by applying the same integral-dimension corollary to the stated polynomial subring.
\end{proof}

\section{Primary decomposition}

In $\Z$, the identity $(12)=(4)\cap(3)$ splits an ideal into two simpler pieces. The ideal $(4)$ is not prime, but if $ab$ is divisible by $4$ and $a$ is not, then $b$ is even. Some power of $b$ is divisible by $4$. This weaker form of primeness is enough for a general decomposition theorem.

Lasker proved the polynomial-ring case in 1905. He is also remembered as world chess champion from 1894 to 1921. Noether showed in 1921 that the ascending chain condition, rather than a special property of polynomials, drives the argument.

\begin{definition}\label{def:primary}
A proper ideal $\mathfrak q\subsetneq R$ is \emph{primary} if, whenever $ab\in\mathfrak q$ and $a\notin\mathfrak q$, some power $b^n$ with $n\ge1$ belongs to $\mathfrak q$.
\end{definition}

Equivalently, every zero divisor in $R/\mathfrak q$ is nilpotent. This follows by taking the images of $a$ and $b$ in the quotient. The radical of a primary ideal records the prime beneath it.

\begin{lemma}\label{lem:primaryradical}
If $\mathfrak q$ is primary, then $\mathfrak p=\sqrt{\mathfrak q}$ is prime. We call $\mathfrak q$ \emph{$\mathfrak p$-primary}.
\end{lemma}

\begin{proof}
The radical is an ideal by the binomial argument in Chapter~4. It is proper because $1^n\notin\mathfrak q$. Suppose $ab\in\sqrt{\mathfrak q}$ and $a\notin\sqrt{\mathfrak q}$. Some $(ab)^m$ lies in $\mathfrak q$, while $a^m$ does not. Applying the primary property gives $(b^m)^k\in\mathfrak q$ for some $k$. Thus $b\in\sqrt{\mathfrak q}$, as required.
\end{proof}

A polynomial example shows why primary pieces need not correspond only to minimal primes. In $K[x,y]$, put $I=(x^2,xy)$. We have
\[
 I=(x)\cap(x^2,y).
\]
Containment from left to right follows on the two generators. Conversely, write an element of the intersection as $xg=ax^2+by$. The prime element $x$ divides $by$ but does not divide $y$, so $b=xc$. Then $xg=ax^2+xyc\in I$. The first component is prime. The second is primary: its quotient is $K[t]/(t^2)$, whose nonunits are multiples of $t$, and those square to zero. Its radical is $(x,y)$, which properly contains $(x)$.

Figure~\ref{fig:v10new-ch05-primary} records the same calculation for monomials. The missing monomial $x$ is why the second component is needed; its square and its product with $y$ both belong to $I$.

\begin{figure}[ht]
\centering
\fig{v10-new-ch05-primary}
\caption{Monomials $x^ay^b$ in $I=(x^2,xy)$ are marked in charcoal. They have $a\ge2$, or $a\ge1$ and $b\ge1$. The rust point $x$ lies in $(x)$ but not in $I$; intersecting with $(x^2,y)$ removes exactly that point from the monomials of $(x)$. The pattern continues upwards and to the right.}
\label{fig:v10new-ch05-primary}
\end{figure}

We first split ideals until they cannot be split further. The ascending chain condition makes that process finite. The key is to show that an unsplittable ideal is primary.

\begin{lemma}\label{lem:irredprimary}
In a Noetherian ring, a proper ideal $I$ that cannot be written as the intersection of two strictly larger ideals is primary.
\end{lemma}

\begin{proof}
For $b\in R$, the \emph{colon ideals}
\[
 (I:b^n)=\{r\in R:rb^n\in I\}
\]
form an increasing chain. Choose $n\ge1$ after this chain has stabilised, so $(I:b^n)=(I:b^{n+1})$. Suppose $ab\in I$ and $a\notin I$. We claim
\[
 I=(I+(a))\cap(I+(b^n)).
\]
Only the reverse containment needs checking. For $w$ in the intersection, write $w=u+ra=v+sb^n$ with $u,v\in I$. It follows that $ra=v-u+sb^n$. Multiply by $b$. Since $ab\in I$, we obtain $sb^{n+1}\in I$. Stabilisation puts $s$ in $(I:b^n)$, so $sb^n\in I$ and $w=v+sb^n\in I$.

The first component $I+(a)$ is strictly larger than $I$. The hypothesis therefore forces $I+(b^n)=I$. Hence $b^n\in I$, proving the primary property. Notice that the forced equality is with a larger ideal; the direction of containment matters here.
\end{proof}

\begin{theorem}[Lasker--Noether]\label{thm:laskernoether}
Every proper ideal in a Noetherian ring is a finite intersection of primary ideals.
\end{theorem}

\begin{proof}
Call a proper ideal \emph{irreducible} if it cannot be split as in the lemma. Suppose there are proper ideals that are not finite intersections of irreducible ideals. By the ascending chain condition in Proposition~\ref{prop:noetherian}, choose a maximal such ideal $I$.

The ideal $I$ is not irreducible, for otherwise its one-term intersection would work. Thus $I=J\cap K$ with both $J$ and $K$ strictly larger. Neither is $R$, since that would make the intersection equal to the other component. Maximality of $I$ implies that each of $J$ and $K$ is a finite intersection of irreducible ideals. Combining these lists gives such a representation for $I$, a contradiction. Every proper ideal is therefore a finite intersection of irreducible ideals. The preceding lemma makes all of these components primary.
\end{proof}

A primary decomposition $I=\mathfrak q_1\cap\cdots\cap\mathfrak q_r$ is \emph{minimal} when its radicals are distinct and no component can be removed. We can always arrange this. Delete redundant components first, and merge components with the same radical. Their intersection is still primary with that radical: if $ab$ belongs to both and $b$ avoids the common radical, primaryness forces $a$ into both. The radical of their intersection is the common radical, since the family is finite.

The primary ideals themselves may depend on the choices. Their radicals do not. Colon ideals let us recover them directly from $I$.

\begin{theorem}\label{thm:assprimes}
Let $I=\mathfrak q_1\cap\cdots\cap\mathfrak q_r$ be a minimal primary decomposition in a Noetherian ring, and put $\mathfrak p_i=\sqrt{\mathfrak q_i}$. A prime ideal is of the form $\sqrt{(I:x)}$ for some $x\in R$ if and only if it is one of the $\mathfrak p_i$. These \emph{associated primes} of $I$ are independent of the decomposition. Their minimal members are exactly the primes minimal over $I$. Every other associated prime contains a minimal one and is called \emph{embedded}.
\end{theorem}

\begin{proof}
Start with a $\mathfrak p$-primary ideal $\mathfrak q$. If $x\in\mathfrak q$, its colon $(\mathfrak q:x)$ is $R$. If $x\notin\mathfrak q$, we claim
\[
 \sqrt{(\mathfrak q:x)}=\mathfrak p.
\]
An element $a\in\mathfrak p$ has some power in $\mathfrak q$, so that power multiplied by $x$ is also in $\mathfrak q$. This proves one inclusion. Conversely, if $a^mx\in\mathfrak q$ and $a\notin\mathfrak p$, primaryness implies $x\in\mathfrak q$: otherwise some power of $a^m$ would lie in $\mathfrak q$. This contradicts the choice of $x$, proving the claim. The same reasoning shows $(\mathfrak q:a)=\mathfrak q$ whenever $a\notin\mathfrak p$.

Colon distributes over intersections, and radical distributes over finite intersections. For the latter assertion, choose a power witnessing membership in each radical, then use the maximum of the finitely many powers. Hence
\[
 \sqrt{(I:x)}=\bigcap_{x\notin\mathfrak q_i}\mathfrak p_i.
\]
The intersection over an empty index set is $R$, which is not prime. Minimality gives, for each $i$, an $x$ belonging to every other component but not to $\mathfrak q_i$. The displayed formula then gives precisely $\mathfrak p_i$.

Conversely, suppose this finite intersection is a prime $\mathfrak p$. It must equal one of its members. Indeed, if it contained none of the $\mathfrak p_i$ as an ideal, choose $a_i\in\mathfrak p_i\setminus\mathfrak p$. Their product lies in every $\mathfrak p_i$ and hence in $\mathfrak p$, contradicting primeness. Thus some $\mathfrak p_i\subseteq\mathfrak p$, while the intersection gives $\mathfrak p\subseteq\mathfrak p_i$.

Finally, every prime $P$ containing $I$ contains one component $\mathfrak q_i$. Otherwise choose an element of each $\mathfrak q_i$ outside $P$ and take their product. Therefore $P$ contains some $\mathfrak p_i$. The primes minimal over $I$ are consequently exactly the minimal members of this finite list. Any member not minimal contains a minimal member, and this is the asserted embedded case.
\end{proof}

In $I=(x^2,xy)$, the associated primes are $(x)$ and $(x,y)$. They can also be found from $(I:y)=(x)$ and $(I:x)=(x,y)$. The second is embedded. In $\Z$, a nonzero prime is already maximal, so associated primes of a nonzero ideal cannot contain one another. The next chapter identifies a larger family of domains with that same property.

\begin{beauty}
The quotient $\Z[x]/(x^2+1)$ is $\Z[i]$. Thus the curve drawn in Figure~\ref{fig:specz} is the spectrum of the Gaussian integers. Over a prime $p\equiv1\pmod4$, $x^2+1$ has two roots and the fibre has two points. For $p\equiv3\pmod4$, it is irreducible and the fibre has one point with residue field of degree $2$. At $2$, it is $(x+1)^2$, so the quotient retains a repeated factor even though its spectrum has only one point. These are the algebraic patterns called splitting, inertia and ramification. They connect polynomial geometry with the arithmetic of primes. Chapter~6 supplies the ideal theory that counts their multiplicities.
\end{beauty}

\section{What we proved, and what comes next}

Localisation keeps the fractions we need and the primes that avoid their denominators. Exactness lets us test a module map locally, and Nakayama turns finite-generation questions in a local ring into linear algebra over its residue field. If you want a concrete model, keep $\Z_{(p)}$ beside the definitions: its units, its maximal ideal and its two prime ideals can all be written down.

For dimension, the useful move is a monic relation. Integral extensions preserve prime-chain lengths; Noether's substitution supplies monic relations; together they reduce dimension to the count for a polynomial ring. Going down is a further theorem for suitable normal domains. We have neither proved nor used it here.

Primary decomposition asks a different question: how does one ideal break into pieces? Existence follows from the ascending chain condition, and colon ideals recover the associated primes despite the possible nonuniqueness of the primary components. Over $\Z$, those pieces are prime powers. Which other domains have this stronger factorisation? Chapter~6 answers with Dedekind domains, where passing to each nonzero prime leaves just a discrete valuation.

\section{Exercises}

\begin{exercise}[$\star$]\label{ex:subringsq}
Let $R\subseteq\Q$ be a unital subring. Prove that $R=S^{-1}\Z$ for a multiplicative set of integers $S$. For a rational prime $p$, determine all ideals of $\Z_{(p)}$ and describe its spectrum.
\end{exercise}

\begin{exercise}[$\star$]\label{ex:golden}
Find an element of the fraction field of $\Z[\sqrt5]$ that is integral over this ring but does not belong to it. Prove that $\Z[(1+\sqrt5)/2]$ is a finitely generated $\Z$-module. Then find a monic polynomial in $\Z[t]$ vanishing at $\sqrt2+\sqrt3$.
\end{exercise}

\begin{exercise}[$\star\star$]\label{ex:localiff}
Prove that a ring is local exactly when its nonunits form an ideal. Show that $R_f=0$ exactly when $f$ is nilpotent. Deduce that the intersection of all prime ideals of $R$ consists precisely of its nilpotent elements.
\end{exercise}

\begin{exercise}[$\star\star$]\label{ex:specrx}
List the prime ideals of $\R[x]$. Construct a bijection between its closed points and $\{z\in\C:\operatorname{Im}z\ge0\}$. Explain separately what happens on the real boundary and above it.
\end{exercise}

\begin{exercise}[$\star\star$]\label{ex:quasicompact}
Show that $\operatorname{Spec}R$ is quasi-compact: every open cover has a finite subcover. First verify that the sets $D(f)=\operatorname{Spec}R\setminus V(f)$ form a basis for its topology.
\end{exercise}

\begin{exercise}[$\star\star$]\label{ex:hyperbola}
For $A=\C[x,y]/(xy-1)$, prove that $A$ is not integral over its subring $\C[x]$, but is integral over $\C[x+y]$. Describe the fibres of the maps $(x,y)\mapsto x$ and $(x,y)\mapsto x+y$ on the complex hyperbola.
\end{exercise}

\begin{exercise}[$\star\star$]\label{ex:v10new-ch05-torsion}
Let $M=\Z/60\Z$, and let $M_{(p)}$ mean localisation at $\Z\setminus(p)$. Compute $M_{(2)}$, $M_{(3)}$ and $M_{(5)}$, including the localisation maps. Prove that the map from $M$ to their direct product is an isomorphism. Compute $M_{(p)}$ for every other rational prime $p$, and explain why each of the three displayed factors is needed to detect all elements of $M$.
\end{exercise}

\begin{exercise}[$\star\star$]\label{ex:v10new-ch05-dual}
Let $K$ be a field and $A=K[\varepsilon]/(\varepsilon^2)$. Determine all its primes, its residue field and its dimension. Show that the quotient map $A\to K$ induces a homeomorphism of spectra, although the rings are not isomorphic. More generally, prove that quotienting any ring by an ideal of nilpotent elements does not change its spectrum as a topological space.
\end{exercise}

\begin{exercise}[$\star\star$]\label{ex:v10new-ch05-cusp}
For a field $K$ and an indeterminate $u$, put $A=K[u^2,u^3]\subset K[u]$. Prove that $A$ is Noetherian, has fraction field $K(u)$ and dimension $1$, and is not integrally closed. Prove that its integral closure in $K(u)$ is $K[u]$. Identify the precise role of the missing element $u$ in each part of the argument.
\end{exercise}

\begin{exercise}[$\star\star\star$]\label{ex:v10new-ch05-embedded}
In $R=K[x,y]$, put $I=(x^2,xy)$. For $c\in K$, set $\mathfrak q_c=(x^2,y-cx)$. Prove that $I=(x)\cap\mathfrak q_c$ is a minimal primary decomposition for every $c$. Show that $\mathfrak q_c\ne\mathfrak q_d$ if $c\ne d$. Compute $(I:x)$ and $(I:y)$, and explain which part of the decomposition stays fixed as $c$ varies.
\end{exercise}

\section{Hints}

\begin{hint}{ex:subringsq}
For a reduced fraction $a/b\in R$, use integers $u,v$ with $ua+vb=1$ to get $1/b\in R$. In $\Z_{(p)}$, separate the power of $p$ from the unit part of a nonzero element.
\end{hint}

\begin{hint}{ex:golden}
The golden ratio satisfies a monic quadratic. For the last task, first isolate $2\sqrt6$ in the square of the proposed root, then square again.
\end{hint}

\begin{hint}{ex:localiff}
A proper ideal contains no unit. For the nilradical assertion, a nonnilpotent $f$ gives a nonzero ring $R_f$; choose a maximal ideal there and contract it.
\end{hint}

\begin{hint}{ex:specrx}
Use that $\R[x]$ is a PID and that a real irreducible polynomial has degree $1$ or $2$. A quadratic with nonreal roots supplies one conjugate pair.
\end{hint}

\begin{hint}{ex:quasicompact}
Refine a cover to basic opens. If the $D(f_\lambda)$ cover, the ideal generated by the $f_\lambda$ is in no maximal ideal. An expression for $1$ uses only finitely many generators.
\end{hint}

\begin{hint}{ex:hyperbola}
Identify $A$ with $\C[x,x^{-1}]$. An integral equation for $x^{-1}$ would become, after multiplication by a power of $x$, a polynomial identity with constant term $1$. For the other subring, use $y^2-(x+y)y+1=0$.
\end{hint}

\begin{hint}{ex:v10new-ch05-torsion}
Use the integer Chinese remainder theorem to split $M$ into its $4$, $3$ and $5$ parts. Inverting $3$ kills any group killed by $3$, and inverting $5$ kills any group killed by $5$.
\end{hint}

\begin{hint}{ex:v10new-ch05-dual}
Every prime contains a nilpotent element. Use the quotient correspondence for ideals and the formula for closed sets $V(I)$.
\end{hint}

\begin{hint}{ex:v10new-ch05-cusp}
The element $u$ satisfies $T^2-u^2=0$ over $A$, and $u=u^3/u^2$ in its fraction field. To rule out further integral elements, use that the UFD $K[u]$ is integrally closed.
\end{hint}

\begin{hint}{ex:v10new-ch05-embedded}
Modulo $\mathfrak q_c$, substitute $y=cx$ and $x^2=0$. For $f=xg$, its image vanishes exactly when the constant coefficient of $g$ is zero.
\end{hint}

\section{Solutions}

\begin{solution}{ex:subringsq}
\idea Admitting a reduced fraction forces its denominator to be invertible. At a fixed prime, only the exponent of that prime matters.
\writeup Put $S=\{b\in\Z_{>0}:1/b\in R\}$. It contains $1$ and is multiplicative. If $a/b\in R$ is reduced with $b>0$, choose $u,v\in\Z$ with $ua+vb=1$. Since the subring is unital, it contains $\Z$, and
\[
 1/b=u(a/b)+v\in R.
\]
Thus every element of $R$ belongs to $S^{-1}\Z$. Conversely, every allowed fraction belongs to $R$ by the definition of $S$.

In $\Z_{(p)}$, a nonzero element has a unique expression $p^ku$ with $k\ge0$ and $u$ a unit. If an ideal is nonzero, choose an element with the smallest possible $k$. Dividing its unit factor shows that $p^k$ is in the ideal, while every other element is a multiple of $p^k$. The ideals are therefore $0$ and $p^k\Z_{(p)}$ for $k\ge0$, including the unit ideal at $k=0$. Among proper ideals, the primes are $0$ and $p\Z_{(p)}$. The latter is the closed point; the former has closure the whole two-point spectrum.
\end{solution}

\begin{solution}{ex:golden}
\idea A monic equation gives integrality; a coefficient check shows the element is missing from the ring.
\writeup Set $\varphi=(1+\sqrt5)/2$. It belongs to $\Q(\sqrt5)$ and satisfies $\varphi^2-\varphi-1=0$. It is therefore integral over $\Z$, and hence over $\Z[\sqrt5]$. Its coefficient of $\sqrt5$ is $1/2$, so it does not belong to that ring. This proves failure of integral closure. Since $\varphi^2=\varphi+1$, reduction of powers gives
$\Z[\varphi]=\Z+\Z\varphi$, a module generated by two elements.

For $t=\sqrt2+\sqrt3$, the successive equalities are $t^2=5+2\sqrt6$ and $(t^2-5)^2=24$. Expanding the last equality gives the monic equation $t^4-10t^2+1=0$.
\end{solution}

\begin{solution}{ex:localiff}
\idea Units cannot lie in a proper ideal, while localisation makes a chosen nonnilpotent element into a unit without killing the ring.
\writeup In a local ring, the nonunits are exactly the unique maximal ideal, as proved above. Conversely, suppose the nonunits form an ideal $\mathfrak m$. It is proper since $1$ is a unit. Every proper ideal consists entirely of nonunits and is therefore contained in $\mathfrak m$. Hence $\mathfrak m$ is the unique maximal ideal. In the zero ring the nonunits form the empty set, which is not an ideal; this agrees with the fact that the ring has no maximal ideal and is not local.

The fraction test gives $R_f=0$ exactly when $1/1=0$, equivalently when $f^n=0$ for some $n\ge0$. In a nonzero ring this exponent must be positive; in the zero ring $f=0$ is nilpotent as well. A nilpotent element belongs to every prime by repeated use of primeness. If $f$ is not nilpotent, $R_f\ne0$ has a maximal ideal. Its contraction is a prime of $R$ avoiding $f$. Thus no such $f$ lies in the intersection of all primes. For the zero ring, the intersection of the empty family of ideals is the whole ring, whose sole element is nilpotent, so the statement still holds.
\end{solution}

\begin{solution}{ex:specrx}
\idea A real polynomial records either one real root or one conjugate pair of nonreal roots.
\writeup The PID $\R[x]$ has prime ideals $(0)$ and those generated by monic irreducibles. The real factorisation theorem gives the irreducibles $x-a$ for $a\in\R$ and $x^2+bx+c$ with $b^2-4c<0$. Every nonzero prime is maximal and hence closed. Send $(x-a)$ to $a$, and send a quadratic ideal to its unique root in the open upper half-plane. The inverse sends $z$ with positive imaginary part to $((x-z)(x-\bar z))$. This gives the required bijection. The zero ideal is the additional generic point. Complex conjugation pairs the nonreal closed points of $\operatorname{Spec}\C[x]$; this describes the closed-point set over $\R$, without identifying the Zariski topology with the usual Euclidean topology of the half-plane.
\end{solution}

\begin{solution}{ex:quasicompact}
\idea A basic open cover is a claim that certain elements generate the unit ideal, and a formula for $1$ is finite.
\writeup For an ideal $I$, an element of $\operatorname{Spec}R$ is outside $V(I)$ precisely when it omits at least one $f\in I$. Therefore
\[
 \operatorname{Spec}R\setminus V(I)=\bigcup_{f\in I}D(f),
\]
proving that these opens form a basis. Refine an arbitrary open cover to a cover by $D(f_\lambda)$. No prime can contain all the $f_\lambda$. The ideal they generate cannot be proper, since any proper ideal lies in a maximal ideal. Thus $1=\sum_{i=1}^kr_if_{\lambda_i}$ for a finite list. If a prime missed all of these basic opens, it would contain all the $f_{\lambda_i}$ and hence $1$. They therefore form a finite subcover. Choosing one original open containing each selected basic open gives a finite subcover of the original cover. The empty spectrum needs only the empty subcover.
\end{solution}

\begin{solution}{ex:hyperbola}
\idea A monic relation makes a projection finite; the reciprocal coordinate provides an obstruction for the other projection.
\writeup Sending $y$ to $x^{-1}$ identifies $A$ with $\C[x,x^{-1}]$, so the copy of $\C[x]$ is indeed a subring. If $y$ were integral over it, an equation
$y^n+p_{n-1}(x)y^{n-1}+\cdots+p_0(x)=0$ would yield
\[
1+p_{n-1}(x)x+\cdots+p_0(x)x^n=0.
\]
The left side is a polynomial with constant term $1$, and cannot vanish in the Laurent polynomial ring. Thus $A$ is not integral over $\C[x]$.

Put $s=x+y$. The relation $xy=1$ gives $y^2-sy+1=0$, so $y$ is integral over $\C[s]$. As $x=s-y$, the ring $A=\C[s,y]$ is integral over $\C[s]$ as well. The element $s$ is transcendental: any nonzero polynomial in $x+x^{-1}$ has a nonzero highest Laurent power of $x$. Hence this is a polynomial subring.

The projection to $x$ has one point above each nonzero complex number and none above $0$. For the projection to $s$, solve $x^2-sx+1=0$ and put $y=1/x$. The roots are always nonzero. There are two distinct points unless $s=2$ or $s=-2$, when there is one repeated root. Every fibre is finite and every complex value of $s$ occurs.
\end{solution}

\begin{solution}{ex:v10new-ch05-torsion}
\idea Localisation keeps the primary part belonging to its prime and kills all parts whose annihilator becomes a unit.
\writeup The Chinese remainder theorem gives $M\cong\Z/4\times\Z/3\times\Z/5$ by reduction. At $(2)$, the integers $3$ and $5$ become units, so the factors killed by them vanish. Every odd integer already acts invertibly on $\Z/4$, leaving $M_{(2)}\cong\Z/4$. Explicitly, $a/s$ maps to $a s^{-1}\pmod4$ for odd $s$, and the original localisation map is reduction modulo $4$. The same reasoning gives $M_{(3)}\cong\Z/3$ and $M_{(5)}\cong\Z/5$, with reduction as the original maps. Their product map is exactly the Chinese remainder isomorphism.

If $p\notin\{2,3,5\}$, the annihilator $60$ becomes a unit, so the localisation is zero. Omitting any one of the three nonzero factors loses information: under the product isomorphism take a tuple nonzero only in the omitted factor. For concrete witnesses, the classes of $15$, $20$ and $12$ are detected only at $2$, $3$ and $5$, respectively. This makes the local zero test visible on a finite module.
\end{solution}

\begin{solution}{ex:v10new-ch05-dual}
\idea Prime ideals cannot distinguish nilpotent elements from zero, but the ring still remembers those elements.
\writeup Every prime of $A$ contains $\varepsilon$, because $\varepsilon^2=0$. The quotient $A/(\varepsilon)=K$ has only the zero prime, so $A$ has just the prime $(\varepsilon)$. It is maximal with residue field $K$, and the spectrum is a one-point space of dimension $0$. The quotient induces the unique homeomorphism from $\operatorname{Spec}K$ to it. Yet $\varepsilon$ is nonzero in $A$ and nilpotent, whereas a field has no nonzero nilpotents. Thus the rings are not isomorphic.

More generally, let $N$ be an ideal all of whose elements are nilpotent. Every prime contains $N$. The quotient correspondence gives a bijection $\operatorname{Spec}(R/N)\to\operatorname{Spec}R$. Under it, $V((I+N)/N)$ corresponds to $V(I)$ for every ideal $I$ of $R$. Closed sets therefore correspond in both directions, proving a homeomorphism. No uniform exponent killing all of $N$ is needed.
\end{solution}

\begin{solution}{ex:v10new-ch05-cusp}
\idea Adding the missing parameter $u$ is an integral extension; the polynomial ring reached in this way is already integrally closed.
\writeup The ring $A$ is generated by two elements over $K$, so Hilbert's basis theorem makes it Noetherian. It is a subring of a domain and has fraction field $K(u)$, since $u=u^3/u^2$. The element $u$ is integral over $A$ by the monic equation $T^2-u^2=0$, but $u\notin A$: the monomials generated by $u^2$ and $u^3$ have exponents $0,2,3,4,\ldots$, never $1$. Thus $A$ is not integrally closed.

As $A[u]=K[u]$, this ring is integral over $A$. The dimension corollary gives $\dim A=\dim K[u]=1$. Every element of $K[u]$ is integral over $A$. Conversely, an element of $K(u)$ integral over $A$ satisfies the same monic equation over $K[u]$, hence belongs to $K[u]$, which is a UFD and integrally closed. This proves the integral-closure assertion in every characteristic. The element $u$ simultaneously recovers the fraction field, witnesses the failure of normality and generates the finite extension that repairs it.
\end{solution}

\begin{solution}{ex:v10new-ch05-embedded}
\idea Change the primary thickening at the embedded point while keeping its intersection with the prime line fixed.
\writeup The quotient by $\mathfrak q_c$ is $K[x]/(x^2)$ via $y=cx$. Its zero divisors are the multiples of $x$, which are nilpotent, so $\mathfrak q_c$ is primary with radical $(x,y)$. Both generators of $I$ vanish in this quotient and are multiples of $x$, giving $I\subseteq(x)\cap\mathfrak q_c$. Conversely, write an element of the intersection as $xg$. Modulo $\mathfrak q_c$, its image is $xg(x,cx)$ modulo $x^2$, equal to the constant coefficient of $g$ times $x$. This vanishes only when that coefficient is zero. Thus $g\in(x,y)$ and $xg\in I$.

The radicals $(x)$ and $(x,y)$ are distinct. Neither component is redundant: $x\in(x)\setminus\mathfrak q_c$, while $y-cx\in\mathfrak q_c\setminus(x)$. The decomposition is minimal. If $c\ne d$, the image of $y-cx$ modulo $\mathfrak q_d$ is $(d-c)x\ne0$, so the primary components differ.

A polynomial $f$ satisfies $fx\in(x^2,xy)$ exactly when $f\in(x,y)$, by cancelling $x$ in the domain. Hence $(I:x)=(x,y)$. Similarly, $fy\in I$ forces $x\mid f$, because $x$ is prime and does not divide $y$, and every such $f$ works; thus $(I:y)=(x)$. The associated primes stay $(x)$ and $(x,y)$ for every $c$. It is the embedded primary component, rather than its radical, that varies.
\end{solution}

\section*{Sources and further reading}
\addcontentsline{toc}{section}{Sources and further reading}

E.~Noether, ``Idealtheorie in Ringbereichen'', \emph{Math.\ Ann.} 83 (1921), 24--66, and ``Der Endlichkeitssatz der Invarianten endlicher linearer Gruppen der Charakteristik $p$'', \emph{Nachr.\ Ges.\ Wiss.\ Göttingen} (1926), 28--35, where
the normalisation lemma appears; W.~Krull, ``Dimensionstheorie in Stellenringen'', \emph{J.\ Reine Angew.\ Math.} 179 (1938), 204--226; I.~S.~Cohen and A.~Seidenberg, ``Prime ideals and integral dependence'', \emph{Bull.\ Amer.\ Math.\ Soc.} 52
(1946), 252--261; E.~Lasker, ``Zur Theorie der Moduln und Ideale'', \emph{Math.\ Ann.} 60 (1905), 20--116. The picture of $\operatorname{Spec}\Z[x]$ is in D.~Mumford, \emph{The Red Book of Varieties and Schemes} (Springer Lecture Notes in Mathematics 1358, 1988), first circulated as mimeographed notes in 1967.

For more: M.~F.~Atiyah and I.~G.~Macdonald, \emph{Introduction to Commutative Algebra} (1969), Chapters~3, 5 and 11; D.~Eisenbud, \emph{Commutative Algebra with a View Toward Algebraic Geometry} (1995), Chapters~2, 4, 8 and 13;
R.~Vakil, \emph{The Rising Sea: Foundations of Algebraic Geometry}, lecture notes revised many times since 2010, Chapters~3 and 4, for the spectrum and many more pictures.
